% ---------------------------------------------------------------------
% 1. Tipografía y colores
% ---------------------------------------------------------------------
\definefontfamily [mainface] [rm] [Pagella]
\definefontfamily [mainface] [ss] [Iwona]
\definefontfamily [mainface] [tt] [ModernMono]
\setupbodyfont     [mainface,11pt]

\definecolor [inkblue]    [r=0.11, g=0.16, b=0.32]
\definecolor [accent]     [r=0.72, g=0.20, b=0.20]
\definecolor [accentb]    [r=0.10, g=0.35, b=0.45]
\definecolor [resultgreen][r=0.13, g=0.42, b=0.24]
\definecolor [softgray]   [r=0.45, g=0.45, b=0.45]
\definecolor [panelbg]    [r=0.96, g=0.96, b=0.94]
\definecolor [resultbg]   [r=0.94, g=0.97, b=0.94]
\definecolor [codebg]     [r=0.965,g=0.965,b=0.98]
\definecolor [codeframe]  [r=0.80, g=0.80, b=0.85]

\setupcolors [state=start]

% ---------------------------------------------------------------------
% 2. Página y márgenes
% ---------------------------------------------------------------------
\setuppapersize   [letter] 
\setuplayout      [
  topspace=2cm, backspace=3cm,
  header=1cm, footer=1cm,
  width=middle, height=middle,
  leftmargin=0cm, rightmargin=0cm,
]
\setupwhitespace [medium]
\setupindenting  [never]

% ---------------------------------------------------------------------
% 3. Encabezados de sección
% ---------------------------------------------------------------------
\definehead   [tutorialsection] [section]

\setuphead [tutorialsection]
  [style=\bf,
   color=inkblue,
   number=yes,
   numbercolor=accent,
   numberstyle=\bf]

% ---------------------------------------------------------------------
% 4. Bloques de código
% ---------------------------------------------------------------------
\definetyping  [contextcode] [
  option=TEX,
  space=on,
  style=\tt,
]

\setupbackgrounds [contextcode] [background=color, backgroundcolor=codebg]

\setuptyping [contextcode] [
  frame=on,
  framecolor=codeframe,
  rulethickness=0.4pt,
  leftmargindistance=1em,
  rightmargindistance=1em,
  topspace=0.6em,
  bottomspace=0.6em,
]

\definehighlight [texcmd] [style=\bf, color=accent]

% ---------------------------------------------------------------------
% 5. Cajas: Objetivo, Acción, Explicación, Observación, Resultado
% ---------------------------------------------------------------------
\defineframedtext
  [objetivo]
  [width=broad, frame=off, background=color, backgroundcolor=panelbg,
   leftframe=on, framecolor=accentb, rulethickness=2pt, offset=1.2em]

\defineframedtext
  [accion]
  [width=broad, frame=off, background=color, backgroundcolor=panelbg,
   leftframe=on, framecolor=accent, rulethickness=2pt, offset=1.2em]

\defineframedtext
  [explicacion]
  [width=broad, frame=off, background=color, backgroundcolor=panelbg,
   leftframe=on, framecolor=accentb, rulethickness=2pt, offset=1.2em]

\defineframedtext
  [observacion]
  [width=broad, frame=off, background=color, backgroundcolor=panelbg,
   leftframe=on, framecolor=accent, rulethickness=2pt, offset=1.2em]

\defineframedtext
  [resultado]
  [width=broad, frame=off, background=color, backgroundcolor=resultbg,
   leftframe=on, framecolor=resultgreen, rulethickness=2pt, offset=1.2em]

\define[1]\Objetivo{\blank[small]\startobjetivo{\bf Objetivo. }#1\stopobjetivo\blank[small]}
\define[1]\Accion{\blank[small]\startaccion{\bf Acción. }#1\stopaccion\blank[small]}
\define[1]\Explicacion{\blank[small]\startexplicacion{\bf Explicación. }#1\stopexplicacion\blank[small]}
\define[1]\Observacion{\blank[small]\startobservacion{\bf Observación. }#1\stopobservacion\blank[small]}
\define[1]\Resultado{\blank[small]\startresultado{\bf\color[resultgreen]{Resultado.} }#1\stopresultado\blank[small]}

% ---------------------------------------------------------------------
% 6. Interacción
% ---------------------------------------------------------------------
\setupinteraction [state=start, color=accentb, contrastcolor=accentb,
                   title={Formato de caracteres en ConTeXt}]

% ---------------------------------------------------------------------
% 7. Pie de página
% ---------------------------------------------------------------------
\setupheadertexts  [][]
\setupfootertexts  [][{\ss\small\color[softgray]{\pagenumber}}]


\environment env_tutorial
% =====================================================================
\starttext

% ------------------------ PORTADA -------------------------------------
\startfrontmatter

\startstandardmakeup[pagestate=start]
  \blank[3*big]
  \startcolor[accent]\hrule height 2pt\stopcolor
  \blank[medium]
  {\ss\bf\switchtobodyfont[24pt]\color[inkblue]{Tutorial}}\par
  \blank[small]
  {\ss\bf\switchtobodyfont[17pt]\color[inkblue]{Domina el formato de}}\par
  {\ss\bf\switchtobodyfont[17pt]\color[accent]{caracteres, palabras y espacios}}\par
  \blank[medium]
  \startcolor[accentb]\hrule height 0.8pt\stopcolor
  \blank[2*big]
  {\it\color[softgray]{Mayúsculas, versalitas, super/subíndices, texto
   literal y espaciado: el detalle tipográfico en ConTeXt.}}
  \vfill
  {\ss\small\color[softgray]{Guía práctica paso a paso}}
\stopstandardmakeup

\stopfrontmatter

% ------------------------ RESUMEN ---------------------------------------
\startfrontmatter

\tutorialsection{Resumen}
En este tutorial aprenderemos a controlar el formato de caracteres,
palabras y el espaciado horizontal en ConTeXt, paso a paso. Al
final, podrás modificar mayúsculas y minúsculas, crear versalitas,
usar superíndices y subíndices, mostrar código literal y gestionar
el espacio entre palabras. Cada sección incluye el código fuente
\emph{y} el resultado real, para que veas el efecto de inmediato.

\blank[medium]
{\ss\bf\color[inkblue]{Contenido}}
\startcolor[codeframe]\hrule\stopcolor
\blank[small]
\placelist[tutorialsection][criterium=all,alternative=c]

\stopfrontmatter

% ------------------------ CUERPO ----------------------------------------
\startbodymatter

\tutorialsection{Prepara tu espacio de trabajo}

Primero, crea una carpeta llamada \tt{mi\_proyecto\_formato} y dentro
de ella, un archivo llamado \tt{documento.tex} con este contenido
inicial:

\blank[small]
\startcontextcode
\starttext

\chapter{Formato de caracteres}

Este es mi documento para aprender a formatear caracteres y
palabras en ConTeXt.

\chapter{Siguiente capítulo}

Más contenido para practicar.

\stoptext
\stopcontextcode
\blank[small]

Compila con \tt{context documento.tex} y observa el resultado.
Obtendrás un PDF con el formato básico. Ahora, vamos a modificarlo.

\tutorialsection{Cambia mayúsculas y minúsculas}

\Objetivo{Convertir texto a mayúsculas o minúsculas de forma
automática.}

\Accion{Usa los comandos \type{\word}, \type{\Word}, \type{\Words} y
\type{\WORD}:}

\blank[small]
\startcontextcode
\starttext

\word        {esTo es UNa PRueba}
\Word        {esTo es UNa PRueba}
\Words       {esTo es UNa PRueba}
\WORD        {esTo es UNa PRueba}
\WORDS       {esTo es UNa PRueba}

\stoptext
\stopcontextcode
\blank[small]

\Resultado{
\startitemize[packed,1]
\item \word{esTo es UNa PRueba}
\item \Word{esTo es UNa PRueba}
\item \Words{esTo es UNa PRueba}
\item \WORD{esTo es UNa PRueba}
\item \WORDS{esTo es UNa PRueba}
\stopitemize}

Observa los resultados:

\startitemize[packed,1]
\item \type{\word}: todo en minúsculas.
\item \type{\Word}: primera letra en mayúscula.
\item \type{\Words}: primera letra de cada palabra en mayúscula.
\item \type{\WORD} y \type{\WORDS}: todo en mayúsculas.
\stopitemize

\Observacion{\type{\WORD} y \type{\WORDS} hacen lo mismo.}

\tutorialsection{Crea versalitas (falsas pequeñas mayúsculas)}

\Objetivo{Aplicar el efecto de pequeñas mayúsculas a un texto.}

\Accion{Usa los comandos \type{\cap} y \type{\Cap}:}

\blank[small]
\startcontextcode
\starttext

\cap{texto en versalitas}
\Cap{Texto en versalitas con primera letra grande}

\stoptext
\stopcontextcode
\blank[small]

\Resultado{\cap{texto en versalitas} \quad
\Cap{Texto en versalitas con primera letra grande}}

Verás el texto en versalitas. \type{\cap} convierte todo en
mayúsculas reducidas, mientras que \type{\Cap} mantiene la primera
letra de cada palabra en tamaño normal.

\Explicacion{Puedes anidar \type{\cap} para reducir aún más el
tamaño: \cap{las personas que han amasado su \cap{capital} a costa
de los demás\ldots}}

\tutorialsection{Usa superíndices y subíndices}

\Objetivo{Colocar texto en superíndice (arriba) o subíndice (abajo).}

\Accion{Usa los comandos \type{\high}, \type{\low} y \type{\lohi}:}

\blank[small]
\startcontextcode
\starttext

Este es un texto \high{arriba} y otro \low{abajo}.
También una combinación \lohi{abajo}{arriba}.

\stoptext
\stopcontextcode
\blank[small]

\Resultado{Este es un texto \high{arriba} y otro \low{abajo}.
También una combinación \lohi{abajo}{arriba}.}

\Observacion{Estos comandos funcionan fuera del modo matemático.
Dentro del modo matemático, usa \type{^} y \type{_}.}

\tutorialsection{Muestra código literal (verbatim)}

\Objetivo{Mostrar texto tal cual, sin que se interpreten comandos.}

\Accion{Usa el comando \type{\type} para texto corto y el entorno
\type{starttyping} para bloques largos:}

\blank[small]
\startcontextcode
\starttext

\type{Texto literal con \comandos y {caracteres} especiales}

\starttyping
Este es un bloque de texto literal.
Los comandos como \starttext no se ejecutan.
Se muestran exactamente como se escriben.
\stoptyping

\stoptext
\stopcontextcode
\blank[small]

Verás el texto en formato monoespaciado (como máquina de escribir)
sin que los comandos se ejecuten.

\Observacion{\type{\type} es para una línea; \type{starttyping} es
para múltiples líneas.}

\tutorialsection{Usa un delimitador personalizado en type}

\Objetivo{Incluir llaves o caracteres especiales dentro de un
\type{\type}.}

\Accion{Usa un carácter delimitador diferente en \type{\type}:}

\blank[small]
\startcontextcode
\starttext

\type {Este texto contiene llaves: { } }
\type |Este texto contiene llaves: { } |
\type *Este texto contiene llaves: { } *

\stoptext
\stopcontextcode
\blank[small]

Todos los ejemplos muestran el texto correctamente. El delimitador
(el primer carácter después de \type{\type}) define dónde empieza y
termina el argumento.

\Explicacion{Si tu texto contiene llaves \type{{ }}, usa un
delimitador diferente como \type{|} o \type{*}.}

\tutorialsection{Añade espacio horizontal entre caracteres}

\Objetivo{Separar caracteres dentro de una palabra para dar un
efecto visual.}

\Accion{Usa el comando \type{\stretched}:}

\blank[small]
\startcontextcode
\starttext

\stretched[factor=0.5, width=5cm]{Texto separado}
\stretched[factor=1.0, width=5cm]{Texto separado}
\stretched[factor=1.5, width=5cm]{Texto separado}

\stoptext
\stopcontextcode
\blank[small]

\Resultado{
\stretched[factor=0.5, width=8cm]{Texto separado}\par
\stretched[factor=1.0, width=8cm]{Texto separado}\par
\stretched[factor=1.5, width=8cm]{Texto separado}}

Verás cómo el espacio entre caracteres aumenta con el factor.
\type{factor} controla la cantidad de separación.

\Observacion{Si solo usas \type{width}, el texto se estira para
ocupar ese ancho.}

\tutorialsection{Añade espacio horizontal entre palabras}

\Objetivo{Controlar el espacio entre palabras en una línea.}

\Accion{Usa comandos como espacio fino, \type{\enskip}, \type{\quad}
y \type{\qquad}:}

\blank[small]
\startcontextcode
\starttext

Espacio normal: Palabra1 Palabra2

Espacio fino: Palabra1\,Palabra2

Medio em: Palabra1\enskip Palabra2

Un em (quad): Palabra1\quad Palabra2

Dos ems (qquad): Palabra1\qquad Palabra2

\stoptext
\stopcontextcode
\blank[small]

\Resultado{
Espacio normal: Palabra1 Palabra2\par
Espacio fino: Palabra1\,Palabra2\par
Medio em: Palabra1\enskip Palabra2\par
Un em (quad): Palabra1\quad Palabra2\par
Dos ems (qquad): Palabra1\qquad Palabra2}

Observa cómo los diferentes comandos crean espacios de distintos
tamaños. El espacio fino es útil para separar miles: 1\,000\,000.

\tutorialsection{Alinea texto con espacio flexible}

\Objetivo{Usar espacio elástico para alinear texto a la derecha o
centrado.}

\Accion{Usa \type{\hfill} para empujar texto:}

\blank[small]
\startcontextcode
\starttext

\hfill Texto a la derecha

Texto a la izquierda \hfill Texto a la derecha

\bf{Centrado} \hfill \bf{Centrado}

\stoptext
\stopcontextcode
\blank[small]

\Resultado{
\hfill Texto a la derecha\par
Texto a la izquierda \hfill Texto a la derecha\par
\bf{Centrado} \hfill \bf{Centrado}}

\type{\hfill} ocupa todo el espacio disponible, empujando el texto
a los bordes.

\Explicacion{Usa \type{\wordright{texto}} para alinear una palabra
a la derecha y saltar de línea:}

\blank[small]
\startcontextcode
\wordright{Palabra a la derecha}
\stopcontextcode
\blank[small]

\tutorialsection{Crea palabras compuestas}

\Objetivo{Indicar a ConTeXt que dos palabras forman una palabra
compuesta para una correcta división (guionización).}

\Accion{Usa el separador \type{||} (dos barras verticales):}

\blank[small]
\startcontextcode
\starttext

Español||argentino

Unir|/|separar

(inter||)||espacio

\stoptext
\stopcontextcode
\blank[small]

\Resultado{Español||argentino \quad Unir|/|separar \quad
(inter||)||espacio}

ConTeXt trata estas secuencias como una sola palabra compuesta,
aplicando reglas de división adecuadas.

\Observacion{El separador por defecto es el guion. Puedes cambiarlo
con \type{\setuphyphenmark}.}

\tutorialsection{Prepara un documento final}

\Objetivo{Usar diferentes técnicas de formato en un documento.}

\Accion{Aplica todos los métodos en un mismo texto:}

\blank[small]
\startcontextcode
\starttext

\chapter{Formatos de texto}

\WORD{texto en mayúsculas}

\cap{Texto en versalitas}

Texto \high{arriba} y \low{abajo}.

Código: \type{print("Hola mundo")}

Espacios: \qquad separación \quad grande.

\stretched[factor=0.5, width=8cm]{Texto separado}

\hfill \wordright{Final}

\stoptext
\stopcontextcode
\blank[small]

Este documento muestra todos los métodos de formato. Puedes usarlo
como referencia.

\tutorialsection{Ejercicio}

Recrea y mejora el siguiente texto:

%\setuppapersize[letter]
%\setupbodyfont[11pt, palatino]
%\setuplayout[backspace=2.5cm, width=16cm, topspace=2cm, height=24cm]
%\setupinterlinespace[small]
%
%% Estilo elegante de títulos
%\setuphead[section][style=\bfb, color=darkred, before={\blank[small]}, after={\blank[small]}]
%\setuphead[subject][style=\bfa, color=darkblue, before={\blank[medium]}, after={\blank[small]}]
%
%\subject{1.~Mayúsculas y minúsculas}
%
%\type{\word} convierte a minúsculas, \type{\Word} pone la inicial en
%mayúscula, \type{\Words} pone cada palabra con inicial mayúscula, y
%\type{\WORD} lo convierte todo a mayúsculas.
%
%\blank[small]
%\startalignment[middle]
%{\tfx \word{LA PLUMA ES MÁS PODEROSA QUE LA ESPADA}}\\
%{\tfx \Word{LA PLUMA ES MÁS PODEROSA QUE LA ESPADA}}\\
%{\tfx \Words{LA PLUMA ES MÁS PODEROSA QUE LA ESPADA}}\\
%{\tfx \WORD{la pluma es más poderosa que la espada}}
%\stopalignment
%
%\blank[small]
%\subject{2.~Versalitas (falsas pequeñas mayúsculas)}
%
%\type{\cap} pasa todo el texto a versalitas; \type{\Cap} convierte
%solo la primera letra en versalita.
%
%\blank[small]
%\startalignment[middle]
%{\tfx \cap{prólogo} \quad \Cap{prólogo} \quad \Cap{cervantes} \quad \cap{el quijote}}
%\stopalignment
%
%\blank[small]
%\subject{3.~Superíndices y subíndices}
%
%\type{\high} coloca un superíndice, \type{\low} un subíndice, y
%\type{\lohi} coloca ambos a la vez.
%
%\blank[small]
%\startalignment[middle]
%{\tfx H\high{2}O \quad CO\low{2} \quad x\lohi{2}{3} \quad E = mc\high{2}}
%\stopalignment
%
%\blank[small]
%\subject{4.~Código literal}
%
%\type{\type} muestra código en línea; el entorno \type{starttyping}
%lo muestra en bloque sin procesar.
%
%\blank[small]
%\starttyping
%\define[1]\saludo{Hola, #1.}
%\saludo{lector}
%\stoptyping
%
%\blank[small]
%\subject{5.~Espaciado entre caracteres y entre palabras}
%
%\type{\stretched} separa las letras (solo aumenta, nunca comprime).
%
%\blank[small]
%\startalignment[middle]
%{\tfx \stretched{ESPACIO} \quad \stretched{ENTRE LETRAS}}
%\stopalignment
%
%\blank[small]
%Entre palabras: \type{\enskip} (medio em), \type{\quad} (un em) y
%\type{\qquad} (dos em).
%
%\blank[small]
%\startalignment[middle]
%{\tfx A\enskip B \quad C\quad D \qquad E\qquad F}
%\stopalignment
%
%\blank[small]
%
%\subject{6.~Espacio elástico y palabras compuestas}
%
%\type{\hfill} reparte el espacio libre de la línea para alinear.
%
%\blank[small]
%\startalignment[middle]
%{\tfx izquierda \hfill centro \hfill derecha}
%\stopalignment
%
%\blank[small]
%Las palabras compuestas se marcan con \type{||}, lo que ayuda a
%\ConTeXt{} a partir correctamente las líneas.
%
%\blank[small]
%\startalignment[middle]
%{\tfx entrada||salida \quad (inter||)espacio \quad intra||palabra}
%\stopalignment

\mainlanguage[es]
\setuppapersize[letter]
\setuplayout[
  topspace=2.0cm,
  backspace=2.5cm,
  width=middle,
  height=middle,
  header=0pt,
  footer=1.5cm,
]

\setupbodyfont[palatino,10.6pt]
\setupinterlinespace[line=1.35em]

% Paleta sobria
\definecolor[acento]     [h=8B2E3C]
\definecolor[gris]       [h=555555]
\definecolor[fondocodigo][h=F4F1EC]

\setuphead[section][
  style=\bfc,number=no,
  color=acento,
  numberstyle=\bfa,
  distance=0.6em,
%  before=1.4em,
%  after=0.5em,
]

\setuptyping[
  style=\tt\tfxx,
  color=gris,
  before={\blank[small]},
  after={\blank[small]},
]

%\starttext

% ─────────────── Encabezado ───────────────
\midaligned{\tfb\sc Afinar la tipografía}
\midaligned{\tfx\color[gris]{Caracteres \cdot\ palabras \cdot\ espacios}}
\blank[big]
\hairline
\blank[medium]

% ─────────────── 1. Caja de casos ───────────────
\section{Mayúsculas y versalitas}

\startalignment[middle]
  \starttabulate[|c|c|c|c|]
  \HL
  \NC \type{\word}  \NC \type{\Word}  \NC \type{\Words}  \NC \type{\WORD}  \NC \NR
  \HL
  \NC \word{tipografía} \NC \Word{tipografía} \NC \Words{tipografía} \NC \WORD{tipografía} \NC \NR
  \HL
  \NC \type{\cap}   \NC \Cap{ConTeXt} \NC \cap{ConTeXt} \NC --- \NC \NR
  \HL
  \stoptabulate
\stopalignment

\blank[small]
El uso de \Cap{versalitas} para acrónimos y de \type{\WORD} para titulares breves
mantiene la \emph{coherencia visual} del documento.

% ─────────────── 2. Sub/superíndices ───────────────
\section{Superíndices y subíndices}

\startalignment[middle]
  \startformula
    E = mc\high{2}, \qquad
    H\low{2}O, \qquad
    a\lohi{n}{k}, \qquad
    \int_{0}^{\infty} e^{-x\high{2}}\,dx = \frac{\sqrt{\pi}}{2}
  \stopformula
\stopalignment

Los tres comandos esenciales son \type{\high}, \type{\low} y \type{\lohi},
útiles tanto en texto corrido (m\high{2}) como en fórmulas.

% ─────────────── 3. Código literal ───────────────
\section{Código literal}

Texto en línea: \type{\stretched{gato}} produce:\thinspace \stretched{gato}. 

Y un bloque completo de código literal se obtiene con el entorno \type{\starttyping}\ldots \type{\stoptyping}:

\starttyping
\def\saludo#1#2{Hola, #1. Bienvenido/a a #2.}
\saludo{Ana}{\ConTeXt}
\stoptyping


% ─────────────── 4. Espaciado ───────────────
%\section{El arte del espacio}
%
%\startalignment[middle]
%  \bf fino:
%  \thinspace|\quad
%  \type{\enskip}:\enskip|\quad
%  \type{\quad}:\quad|\quad
%  \type{\qquad}:\qquad|
%\stopalignment
%
%\blank[small]
%\startalignment[middle]
%  \stretched{Letras separadas con gracia}
%\stopalignment

\blank[medium]
% ─────────────── 5. Alineación con \hfill ───────────────
\section{Alineación elástica}

\startframedtext[width=broad,frame=off,background=color,backgroundcolor=fondocodigo]
  \tfx
  Autor: \hfill \color[gris]{\ConTeXt\ User Group} \crlf
  Fecha: \hfill \color[gris]{2026} \crlf
  Versión: \hfill \color[gris]{1.0}
\stopframedtext

% ─────────────── 6. Palabras compuestas ───────────────
\section{Palabras compuestas}

Cuando escribes \type{Lógico||matemática}, \ConTeXt\ entiende que
\emph{Lógico||matemática} debe permanecer unida, evitando cortes
incómodos al final de línea.

% ─────────────── 7. Configuración avanzada ───────────────
\section{Otras opciones}

\startitemize[packed,1]

\item  Por defecto, \type{\cap} produce \emph{falsas} versalitas.
      Con \type{\setupcapitals[sc=yes]} forzamos el uso de
      \emph{verdaderas} versalitas si la fuente las incorpora.
      Pero también podemos definir variantes con nombre:

\starttyping
\setupcapitals[sc=yes,title=no]
\definecapitals[miscaps][sc=yes]
\stoptyping

      \blank[small]
      \startalignment
        \tfx
        \type{\cap} por defecto:\quad \cap{conTeXt} \quad$\cdot$\quad
        \type{\Cap}:\quad \Cap{ConTeXt}
      \stopalignment

\item {Usar \type{\definecharacterkerning}}: crear
      configuraciones personalizadas para \type{\stretched}.

      \type{\stretched} usa internamente el mecanismo de kerning
      de caracteres. Podemos definir nuestros propios factores con
      \type{\definecharacterkerning} y aplicarlos con
      \type{\setcharacterkerning}:

\starttyping
\definecharacterkerning[muyancho][factor=0.35]
\definecharacterkerning[estrecho][factor=.05]
\setupcharacterkerning[letterspacing][
  style=\spaceskip\dimexpression{
    3.0\spaceamount*\characterkerningparameter{factor}+\spaceamount
  },
]
\stoptyping

\definecharacterkerning[muyancho][factor=0.35]
\definecharacterkerning[estrecho][factor=.05]
\setupcharacterkerning[letterspacing][
  style=\spaceskip\dimexpression{
    3.0\spaceamount*\characterkerningparameter{factor}+\spaceamount
  },
]

      \blank[small]
      \startalignment
        \tfx
        Normal:\quad tipografía \quad$\cdot$\quad
        Estrecho:\quad
          {\setcharacterkerning[estrecho]tipografía} \quad$\cdot$\quad
        Muy ancho:\quad
          {\setcharacterkerning[muyancho]tipografía}\quad$\cdot$\quad
      \type{letterspacing}:\quad{\letterspacing{tipografía}} 
      \stopalignment


\item {Combinar estilos}.
      Los comandos de mayúsculas, versalitas y espaciado pueden
      anidarse con \type{\bf}, \type{\em}, \type{\it} y
      \type{\sc} sin conflicto:

      \blank[small]
      \startalignment[middle]
        \tfx
        \word{conTeXt} \quad
        \bf\Word{conTeXt} \quad
        \em\Words{conTeXt} \quad
        \bf\em\WORD{conTeXt} \quad
        \cap{conTeXt} \quad
        \bf\Cap{conTeXt}
      \stopalignment

      \blank[small]
      \starttyping
\bf\WORD{conTeXt}    % negrita + todo mayúsculas
\em\cap{conTeXt}     % cursiva + versalitas
      \stoptyping

\stopitemize

\definecharacterkerning
  [Title]
  [factor=0.20]

\definefont
  [TitleFont]
  [Serif*default at 18pt]


\startalignment[middle]

{\TitleFont TÍTULO SIN ESTIRAR}

\blank[big]
{\TitleFont\setcharacterkerning[Title]TÍTULO ESTIRADO}

\stopalignment

\blank[big]
\hairline
\blank[small]

\midaligned{
  \tfxx\color[gris]{
    \cap{El buen diseño no se nota: se siente.}
  }
}


\stopbodymatter

% ------------------------ APÉNDICES --------------------------------------
\startappendices

\tutorialsection{¡Has logrado!}

Dominar el formato de caracteres, palabras y espacios en ConTeXt.
Ahora sabes:

\startitemize[packed,1]
\item \bf{Cambiar mayúsculas/minúsculas}: usar \type{\word},
      \type{\Word}, \type{\Words}, \type{\WORD}.
\item \bf{Crear versalitas}: usar \type{\cap} y \type{\Cap} para
      falsas pequeñas mayúsculas.
\item \bf{Superíndices y subíndices}: usar \type{\high}, \type{\low}
      y \type{\lohi}.
\item \bf{Mostrar código literal}: usar \type{\type} y el entorno
      \type{starttyping} para texto sin procesar.
\item \bf{Espaciado entre caracteres}: usar \type{\stretched} para
      separar letras.
\item \bf{Espaciado entre palabras}: usar espacio fino,
      \type{\enskip}, \type{\quad}, \type{\qquad}.
\item \bf{Espacio elástico}: usar \type{\hfill} para alinear texto.
\item \bf{Palabras compuestas}: usar \type{||} para indicar palabras
      compuestas.
\stopitemize

Este flujo de trabajo te permite controlar el formato de texto a
nivel de carácter y palabra, mejorando la legibilidad y el diseño
de tus documentos.

\stopappendices

% ------------------------ BACKMATTER --------------------------------------
\startbackmatter

\tutorialsection{Próximos pasos}

\startitemize[packed,1]
\item \bf{Explora \type{\setupcapitals}}: personaliza el
      comportamiento de \type{\cap} y \type{\Cap}.
\item \bf{Usa \type{\definecharacterkerning}}: crea configuraciones
      personalizadas para \type{\stretched}.
\item \bf{Combina con otros estilos}: aplica negrita o cursiva junto
      con estos comandos.
\item \bf{Recuerda}: el formato excesivo puede dificultar la
      lectura. Úsalo con moderación.
\stopitemize

\stopbackmatter

\stoptext